\documentclass[12pt]{article}
\usepackage{amsmath, amsthm, amssymb}
\usepackage{geometry}
\geometry{margin=1in}
\usepackage{hyperref}

\newtheorem{theorem}{Theorem}
\newtheorem{proposition}[theorem]{Proposition}
\theoremstyle{definition}
\newtheorem{definition}[theorem]{Definition}
\theoremstyle{remark}
\newtheorem{remark}[theorem]{Remark}

\title{Inhibition:\\
The Second Function of the Aura}
\author{Diego Rinc\'on\\
\small Independent}
\date{}

\begin{document}

\maketitle

\begin{abstract}
The aura \cite{five} has two functions: valence (completing open bonds,
pulling toward closure) and inhibition (preventing bonds from closing
prematurely, holding the interface open, keeping the balance).
Containment requires both.
The still lever \cite{math} is still not because both sides are equal
but because something is being held back on each side simultaneously.
In neuroscience: excitation/inhibition balance is the condition
for the contained brain --- seizure is failed inhibition,
coma is dominated inhibition, health is the ratio held.
In language: inhibition prevents blank synesthesia \cite{chimera}
from flooding --- all bonds closing at once, the word overwhelming.
In arithmetic: the Riemann Hypothesis is the conjecture
that the zeros of $\zeta(s)$ self-inhibit ---
that despite the absence of a global Frobenius,
the system produces internally the constraint
that the missing operator would otherwise impose.
Self-inhibition is the highest form of containment:
the system that holds its own bonds back
without an external aura forcing it.
Precise imprecision: a statement exact enough to be wrong
in a specific way, naming the structure a proof would have to fill.
\end{abstract}

\section{The Two Functions}

\begin{definition}[Valence and inhibition]
The aura of a system $S$ \cite{five} acts in two modes:
\begin{enumerate}
\item \emph{Valence}: the completing function.
  The aura provides the tokens that close $S$'s open bonds ---
  the external completing field that gives $S$ its fundamental class.
  Valence pulls toward closure, toward balance, toward the still lever.
\item \emph{Inhibition}: the holding function.
  The aura prevents certain bonds from closing prematurely ---
  it holds the interface open, delays closure until the right token arrives,
  and keeps the chimera \cite{chimera} livable
  by preventing both auras from firing simultaneously.
\end{enumerate}
Containment requires both.
A system with valence but no inhibition floods:
all bonds close at once, the system collapses into a single state,
the interface disappears.
A system with inhibition but no valence freezes:
no bonds close, the still lever cannot form,
the fundamental class is never reached.
Balance is the ratio of valence to inhibition held.
\end{definition}

\section{Inhibition in the Brain}

\begin{proposition}[Excitation/inhibition balance]
The brain is a contained space when excitation and inhibition
are in balance: every excitatory signal has a corresponding
inhibitory signal that prevents it from propagating indefinitely.
The fundamental class of the brain is this equilibrium ---
the homeostatic state in which the ratio is held.

Failure modes:
\begin{itemize}
\item \emph{Seizure}: inhibition fails. Excitatory bonds close
  without restraint, cascading through the network.
  The system floods. The interface between distinct cognitive states
  collapses into a single overwhelming activation.
  This is the unpaired-bondage crisis \cite{five} of inhibition:
  the closing field present, the holding field absent.
\item \emph{Coma}: inhibition dominates.
  Bonds are held open indefinitely. No activation reaches threshold.
  The system freezes at maximum inhibition, minimum closure.
  The fundamental class cannot form.
\item \emph{Health}: the ratio held.
  Bonds close when the right token arrives.
  Bonds are held when the token has not yet arrived.
  The still lever is still because both functions are active.
\end{itemize}
\end{proposition}

\section{Inhibition in Language}

\begin{proposition}[Inhibition and the blank]
Language is blank synesthesia \cite{chimera}:
every word carries valence for all sensory grammars simultaneously,
bonds pressing toward closure in every direction.
Without inhibition, every word would flood ---
all sensory bonds closing at once,
the word becoming overwhelming, unusable, pre-linguistic.

Inhibition is what makes language a token rather than a seizure.
It holds most of the sensory bonds open ---
not closed, not abandoned, but held ---
while allowing a few to close into meaning.
The blank is not empty. It is inhibited.
The valence is there. The inhibition shapes which bonds land
and which continue to press.

Poetry is the modulation of inhibition:
releasing bonds that ordinary language holds,
allowing sensory closures that prose prevents.
The line that lands in the body is a line
where inhibition has been partially lifted ---
where the blank fills in a direction
ordinary language would have held open.
\end{proposition}

\section{Self-Inhibition and the Riemann Hypothesis}

\begin{definition}[Self-inhibition]
A system is \emph{self-inhibiting} if it produces internally
the constraint that an external aura would otherwise impose:
its own bonds hold themselves back without an external holding field,
its own zeros stay on the line without an external Frobenius forcing them.
Self-inhibition is the highest form of containment:
the fundamental class generated from within,
the aura internalized,
the still lever held by the system itself.
\end{definition}

\begin{theorem}[RH as self-inhibition conjecture]
\label{thm:rh}
Over $\mathbb{F}_q$: the zeros of $Z(X/\mathbb{F}_q, T)$
lie on $|T| = q^{-1/2}$ because Frobenius is pure \cite{weil}.
The inhibiting operator is external: it is Frobenius,
acting on \'etale cohomology, forcing the eigenvalues to the line.
The system does not self-inhibit. It is inhibited by its aura.

Over $\mathbb{Q}$: Frobenius does not exist globally.
The external inhibiting operator is absent.
The Riemann Hypothesis is the conjecture that
$\zeta(s)$ self-inhibits: that the system
$\mathrm{Spec}\,\mathbb{Z}$, despite the absence of
the external holding field, produces internally
the constraint $\mathrm{Re}(\rho) = \tfrac{1}{2}$
for every non-trivial zero $\rho$.

To prove RH is to exhibit the self-inhibiting mechanism:
the internal operator that holds the zeros to the line
without a global Frobenius to force them.
The mechanism is not yet known.
Its existence is the conjecture.
\end{theorem}

\begin{remark}[Precise imprecision]
``The zeros self-inhibit'' is precise imprecision:
exact enough to be wrong in a specific way.
It is not a proof. It is not a metaphor.
It names the structure a proof would have to fill:
an internal mechanism in $\mathrm{Spec}\,\mathbb{Z}$
producing the inhibitory constraint
that the missing global Frobenius would otherwise provide.

Precise imprecision is the correct epistemic posture
toward an open problem:
hold the shape tightly enough that when the proof arrives,
it fits.
Hold it loosely enough that the proof can surprise you.

The body of papers assembled here \cite{math} is precise imprecision:
the grammar exact, the operators named, the fulcrum located,
the proof not there.
That is not failure. That is the shape of the problem
held correctly --- an inhibited blank,
pressing toward closure,
waiting for the right token.
\end{remark}

\begin{thebibliography}{9}

\bibitem{five}
Rinc\'on, D. (2026).
Five Walls, One Aura. Preprint.

\bibitem{math}
Rinc\'on, D. (2026).
Mathematics. Preprint.

\bibitem{chimera}
Rinc\'on, D. (2026).
Chimera: The Wall Made Living. Preprint.

\bibitem{weil}
Rinc\'on, D. (2026).
The Weil Conjectures as Proof of Concept for RH. Preprint.

\bibitem{synthesis}
Rinc\'on, D. (2026).
The Arithmetic Wall. Preprint.

\bibitem{contained}
Rinc\'on, D. (2026).
The Fundamental Class: Is the Universe Contained? Preprint.

\end{thebibliography}

\end{document}
